\subsection{From Enterprise to Nation: The Macro-Political Economy of AI Governance}

The enterprise-level contrast documented above points to a deeper question: the choice between open and closed source AI is not an isolated market decision, but is embedded in national-level AI governance strategies. The different judgments made by Chinese and American top leadership about AI's development direction constitute the macro-political conditions for enterprise AI access models.

\subsubsection{The Chinese Path: New Quality Productive Forces and Open-Source Universalism}

Since 2023, China's top leadership has explicitly positioned AI within the strategic framework of \textbf{new quality productive forces}. In multiple public addresses, Xi Jinping has emphasized that AI technologies should serve the transformation of the real economy and high-quality development. The 2024 Government Work Report launched the \textbf{AI+} action initiative, whose core objective is not developing a single cutting-edge model, but diffusing AI capabilities into the capillaries of manufacturing, agriculture, healthcare, and education.

Within this strategic direction, the open-source AI ecosystem has received clear policy support:

\begin{itemize}
\item \textbf{Industrial logic}: if AI is the infrastructure of new quality productive forces, its access threshold must be sufficiently low. Open-source models---such as GLM, DeepSeek, and open-weight versions of Tongyi Qianwen---enable small and medium-sized enterprises and county-level economies to use AI without prohibitive API costs. This reflects the priority of \textbf{productivity logic} over \textbf{business model}: at the national strategic level, what matters is the speed and penetration rate of AI diffusion, not any single company's revenue.
\item \textbf{Digital sovereignty logic}: locally deployed open-source models allow Chinese enterprises to use AI without dependency on American companies' APIs. In the context of US-China technology competition, this is simultaneously an economic choice and a supply-chain security imperative.
\item \textbf{Cognitive universalism logic}: consistent with this paper's cognitive democratization thesis, the recurring policy discourse in Chinese documents about AI empowering SMEs and grassroots governance is essentially about promoting broad spectrum ascent---enabling more people to move from Levels 0--2 to Levels 3--4.
\end{itemize}

Multiple reports from the Development Research Center of the State Council and the China Academy of Information and Communications Technology identify open-source AI as a key pathway toward AI universalism. This stands in systematic contrast to the US AI governance model centered on safety control and export restrictions.

\subsubsection{The American Path: Safety First and Technology Containment}

US AI policy exhibits an internal tension. On one hand, Silicon Valley's innovation tradition values openness and competition; on the other, since 2023, AI safety and alignment have become core themes of bipartisan consensus in Washington.

The Biden Administration's October 2023 Executive Order on Safe, Secure, and Trustworthy AI required the most powerful AI models to undergo safety testing and report results before release. This order objectively strengthened the advantage of large closed-source models---only resource-rich companies (OpenAI, Google, Anthropic) can afford the compliance costs. Distributors of open-source models (such as Meta's Llama series) face a more ambiguous regulatory boundary.

The Trump camp has favored large-scale AI deregulation and emphasized comprehensive competition with China in AI, while retaining the core tool of technology export controls---particularly restrictions on advanced AI chip exports to China. These restrictions, while targeting hardware rather than software, indirectly affect the open-source AI ecosystem: they constrain the hardware resources available to Chinese open-source models during training, partially suppressing the competitiveness of the open-source ecosystem.

Research from the Brookings Institution and the RAND Corporation identifies the core anxiety of US AI governance as the \textbf{safety-competition balance}: simultaneously concerned that AI will be exploited by malicious actors and that excessive regulation will cede innovation to China. This anxiety, in practice, tends toward \textbf{stronger controls}---whether through safety testing requirements or chip export restrictions.

\subsubsection{Two Governance Philosophies Mapped onto Our Framework}

\begin{center}
\begin{tabular}{c|c|c}
\hline
Dimension & Chinese Path & American Path \\
\hline
Core narrative & New productive forces, AI universalism & AI safety, technology competition \\
Governance tools & Industrial policy + open-source support & Safety executive orders + export controls \\
Enterprise AI access & Disintermediation (local deployment) & Re-intermediation (API licensing + PE auditing) \\
Spectrum logic & Lowering spectrum ascent barriers & Reinforcing gatekeeper structures \\
Attitude toward open source & Policy encouragement + industrial push & Safety concerns + regulatory ambiguity \\
\hline
\end{tabular}
\end{center}

A crucial caveat: this is not a simple binary. China has AI safety regulations (e.g., generative AI service management measures), and the US has an active open-source AI community (e.g., the commercial ecosystems around Llama and Mistral). But the \textbf{gravitational center} of the two countries' top-level strategy differs significantly: between diffusion and control, China's policy balance tilts toward diffusion; between openness and safety, America's tilts toward safety.

\subsubsection{Integration with Our Theoretical Framework}

This macro-political contrast reinforces our core arguments:

\begin{enumerate}
\item \textbf{AI access models are political choices, not purely market outcomes}. The choice between open and closed source depends not only on enterprise-level cost-benefit analysis, but on whether national strategy treats AI as \textit{public infrastructure} (requiring disintermediated diffusion) or as a \textit{strategic commodity} (requiring control and licensing).
\item \textbf{The institutional conditions of cognitive democratization}: China's open-source AI policy objectively lowers spectrum ascent barriers---not out of abstract democratic ideals, but out of the practical need for productivity diffusion. This suggests an important insight: the driving force for cognitive democratization need not come from democratic values; it can also come from \textit{developmentalist pragmatism}.
\item \textbf{Open source as a geotechnological choice}: in the context of US-China technology competition, open source is simultaneously a tool for cognitive democratization and a safeguard for technological sovereignty. An AI ecosystem dependent on closed-source APIs faces greater supply disruption risk in times of geopolitical tension.
\end{enumerate}

This analysis also adds an international dimension to our policy implications: cognitive democratization is not only a question of domestic equality, but also a question of power within the international system. When AI becomes the infrastructure of a new era, those who control the entry points to that infrastructure control the cognitive evolution pathways of its users.
